\documentclass[11pt]{article}

\usepackage[margin=1in]{geometry}
\usepackage{graphicx}
\usepackage{booktabs}
\usepackage{amsmath}
\usepackage{amssymb}
\usepackage{caption}
\usepackage{subcaption}
\usepackage{url}
\usepackage[numbers]{natbib}
\usepackage{xcolor}
\usepackage{hyperref}
\hypersetup{
    colorlinks=true,
    linkcolor=black,
    citecolor=blue!60!black,
    urlcolor=blue!70!black,
}
\usepackage{titling}

\title{\textbf{TerrainGen: Generating Playable 2D Platformer Levels\\
with a BFS-Repaired Conditional VAE}\\
\normalsize{AIPI 540 Individual Project --- Technical Report}}

\author{Arnav Mahale\\
\texttt{arnav.mahale@duke.edu}\\
Duke University}

\date{April 2026}

\begin{document}

\maketitle

\begin{abstract}
I study the task of generating short, playable 2D platformer screens in the
style of \emph{SuperTux}. Three generative models, a naive ground-height
baseline, a bigram transition model, and a conditional convolutional VAE,
are trained on a dataset of $\sim$1{,}500 twenty-by-forty tile screens
extracted from the SuperTux open-source game. All three produce outputs
that are visually plausible but almost never playable in isolation
(\,$\approx$0\% BFS reachability). I show that a lightweight
physics-aware repair post-pass (BFS playability search plus greedy
platform-bridging and hazard-carving) lifts playability to
\textbf{84.5\%} for the bigram and \textbf{78.5\%} for the VAE at
medium difficulty, without meaningfully distorting the learned tile
distribution (\(\Delta\text{JS} < 0.005\) across every cell). The
repair ablation is the project's central experiment; its headline
is that for short-horizon tile-grid level generation, the ``hard''
playability constraint is better handled as an explicit
post-processing stage than implicit in the learned distribution. The
system is deployed as a playable web app at
\url{https://arnav161-genterrain.hf.space}.
\end{abstract}

% ---------------------------------------------------------------------------
\section{Problem Statement}

Procedural content generation (PCG) for 2D platformers sits at the
intersection of generative modeling and game-design constraints: the
output must look like a level, be \emph{solvable} (a player can reach
the right edge from the left), and, for a shipped product, be
diverse enough that repeated plays feel fresh. Naive application of
modern generative models to tile grids typically yields outputs that
score well on visual/distributional metrics (e.g., FID-analogues,
Jensen-Shannon divergence) while failing the much-simpler question
\emph{can the player actually finish this level?} This project asks a
concrete engineering question:

\begin{quote}
\emph{If we train a generative model on $\sim$1{,}500 short SuperTux
screens, can we produce an interactive, browser-playable experience
where every generated level is beatable, without requiring the
generative model itself to learn the hard physics constraints?}
\end{quote}

My answer is yes, and the load-bearing piece is not the model, it
is a physics-aware repair pass.

% ---------------------------------------------------------------------------
\section{Data Sources}

All training data derives from the \textbf{SuperTux} open-source 2D
platformer\footnote{\url{https://github.com/SuperTux/supertux}},
MIT-licensed. The raw inputs are:

\begin{itemize}
    \item \texttt{data/raw/tiles.strf} --- the SuperTux tileset
    definition (9{,}296 lines of S-expressions), giving for each tile
    ID its attribute bitfield (SOLID, UNISOLID, SLOPE, HURTS, WATER,
    etc.).
    \item \texttt{data/raw/all\_levels.txt} --- 2{,}433 non-overlapping
    20$\times$40 tile-ID chunks, sliced from SuperTux's $\sim$110
    built-in \texttt{.stl} level files.
\end{itemize}

\paragraph{Tile categorization.}
Raw SuperTux tile IDs encode visual theme (snow, forest, castle) in
addition to gameplay role. Since PCG should be theme-agnostic, I
collapse 6{,}180 unique tile IDs to 8 gameplay categories:
\emph{empty, solid, slope, platform (one-way), bonus, water, hazard,
decoration}. The mapping is built by parsing \texttt{tiles.strf} for
both bulk tile arrays and per-tile boolean attributes. Coverage:
1{,}525 of 1{,}569 non-zero tile IDs observed in levels (\textbf{99.9\%});
the remaining 0.1\% default to \emph{decoration}.

\paragraph{Filtering and splitting.}
Because SuperTux levels can be hundreds of columns tall, many sliced
chunks are mostly sky. I drop any chunk with $<$10\% non-air tiles
(which removes 924 chunks, including all 248 completely-empty ones),
leaving 1{,}509 level chunks. These are split 80/10/10 into
train/val/test with fixed seed 42.

\paragraph{Category imbalance.}
The post-filter category distribution is sharply skewed: empty
(49.9\%), solid (26.9\%), decoration (22.3\%), and everything else
$<$1\% combined. This is load-bearing for training choices (I use
\emph{sqrt}-inverse class weights, not inverse-frequency, see
Section~\ref{sec:models}) and for how the evaluation numbers should
be read (tile-level distribution metrics will be dominated by the
common categories).

% ---------------------------------------------------------------------------
\section{Related Work}

The problem I address has a small but well-developed literature.
\citet{summerville2018procedural} survey PCG via machine learning and
identify \emph{playability guarantees} as a recurring weakness of
purely generative approaches, a framing I adopt directly in the
experiment design. The canonical Super Mario variant of this task is
\citet{volz2018evolving}, which trains a GAN on Mario Bros. screens
and composes them with evolutionary search to enforce a design
objective. My approach differs in two ways: (i) I use a conditional
VAE \citep{kingma2014auto,sohn2015learning} rather than a GAN for
tractability on a small dataset, and (ii) I make the playability
constraint fully deterministic via BFS post-processing, analogous to
\citet{cooper2013paris} and the general ``generate-then-test'' PCG
pattern \citep{togelius2011search}, instead of coupling it into the
search loop.

My bigram classical baseline builds on the long history of grammar-
and transition-based PCG (\citet{smith2010tanagra} among others); my
naive baseline is intentionally minimal, following the course's
insistence on a \emph{truly} uninformative control.

Finally, the \emph{post-hoc repair} approach I ultimately rely on
closely parallels the ``feasibility repair operators'' in search-based
PCG \citep{togelius2011search} and the scene-graph correction stage in
\citet{earle2021illuminating}, though I apply it as a deterministic
post-pass rather than inside a search loop.

% ---------------------------------------------------------------------------
\section{Evaluation Strategy and Metrics}
\label{sec:eval}

My evaluation treats generation as a constrained design problem with
\emph{three} non-negotiables: visual plausibility, playability, and
diversity. Metric choices reflect those objectives.

\paragraph{Playability rate (primary).}
I run a physics-aware breadth-first search from every valid standing
position in the leftmost column; the level passes if any reachable
cell lies in the rightmost column. The BFS is \emph{arc-aware}: jump
moves are rejected unless the mid-flight rectangle
$\{(r,c): r\in[r_\text{peak}, r_\text{max}],\, c\in[c_\text{lo}, c_\text{hi}]\}$
contains no hazard AND both the take-off column (rows
$[n_r, r-1]$ in column $c$) and the two cells above the landing point
are clear of walkable tiles. This extra check mirrors the actual
browser game's physics (\texttt{JUMP\_FORCE=-13}, \texttt{GRAVITY=0.6},
yielding a $\sim$4-tile peak). Without it, BFS over-approximates and
gives false playable labels for levels that are unreachable in the
browser; I verified this myself by repeatedly failing ``playable''
levels live. The constant \texttt{GAME\_JUMP\_PEAK=4} bridges the BFS
abstraction to the continuous parabola.

\paragraph{Jensen-Shannon divergence.}
I compare generated and training tile-category histograms:
\[
\operatorname{JS}(P \| Q) = \tfrac{1}{2} \operatorname{KL}(P \| M) + \tfrac{1}{2} \operatorname{KL}(Q \| M),
\quad M = \tfrac{1}{2}(P+Q).
\]
Lower is better. JS is symmetric and bounded in $[0,\ln 2]$, making
cross-model comparison meaningful. I pre-compute the train-set
histogram once and ship it with the app to avoid loading the full
150MB training JSONL at inference.

\paragraph{Structural metrics.}
\emph{Ground coverage} = fraction of walkable tiles.
\emph{Ground contiguity} = fraction of ground tiles adjacent to another
ground tile (proxy for ``are your platforms archipelagos or islands?'').
\emph{Hazard density} = fraction of hazard tiles. These correlate with
visual legibility more than playability, but diverging from the real
distribution on any of them flags a mode collapse.

\paragraph{Pairwise diversity.}
Mean Hamming distance across 100 sampled level pairs. Higher is more
diverse. I log this to catch the opposite failure mode: a model that
is perfectly playable because it collapses to a single archetype.

\paragraph{Why not classification metrics?}
This is a generative task; precision/recall/F1 are not well-defined.
Tile-level ``accuracy'' (comparing a generated tile to a reference
tile at the same position) is meaningless because there is no ground
truth per position as any position has many legitimate completions.
My closest analogue to a confusion matrix is
Figure~\ref{fig:confusion} (Section~\ref{sec:results}), which shows
each model's playable/unplayable outcome matrix before and after
repair.

% ---------------------------------------------------------------------------
\section{Modeling Approach}
\label{sec:models}

\subsection{Data processing pipeline}
Each stage has a distinct role:
\begin{enumerate}
    \item \textbf{Parse tileset} (\texttt{build\_tile\_mapping.py}):
    map 6{,}180 tile IDs to 8 gameplay categories from the SuperTux
    attribute bitfield. \emph{Rationale:} theme-invariant targets cut
    vocabulary by $\sim$800$\times$ and make the tiny dataset tractable.
    \item \textbf{Build features} (\texttt{build\_features.py}):
    remap the 2{,}433 raw chunks and drop those with $<$10\% non-air
    content. \emph{Rationale:} near-empty chunks would bias the model
    toward outputting sky.
    \item \textbf{Assign difficulty buckets} (\texttt{difficulty.py}):
    score each chunk by a structural difficulty (gaps,
    hazards, height variance) and bin into easy / medium /
    hard. \emph{Rationale:} gives the VAE a conditioning signal I
    can expose via a UI slider.
    \item \textbf{Bake layout into training data}
    (\texttt{repair.enforce\_layout}): force top 4 rows to empty sky
    and bottom 2 rows to solid floor before training the VAE.
    \emph{Rationale:} these bands are enforced at inference too, so
    the VAE should not waste capacity modeling them.
\end{enumerate}

\subsection{Hyperparameter tuning strategy}
I chose three configurations informed by failure modes
observed on the train/val split:
\begin{itemize}
    \item Inverse-frequency class weights collapsed the VAE to near-
    uniform output (rare hazards over-represented). I switched to
    \emph{sqrt}-inverse weights
    $w_c \propto 1/\sqrt{\operatorname{freq}(c)}$, capped at 1.5, which
    recovers a realistic histogram.
    \item KL weight $\beta = 0.5$ drove posterior collapse; I use
    $\beta = 0.1$.
    \item Latent dimension 64 (vs.\ the prior project's 32) gave
    better reconstruction on val without harming diversity.
    \item Epochs 150, early selection on val cross-entropy. No
    learning-rate search; Adam at $10^{-3}$ with
    \texttt{ReduceLROnPlateau} was sufficient.
    \item Classifier-free guidance: drop conditioning with probability
    $p=0.15$ during training; at inference use guidance scale 2.0.
\end{itemize}

\subsection{Models evaluated}

\paragraph{Naive baseline (\texttt{model\_naive.py}).}
Learn the mean and standard deviation of per-column ground height from
training data ($\bar{h} \approx 6.9$, $\sigma_h \approx 8.1$). At
generation time, sample a per-column height and fill solid below.
\emph{Why this?} Strictly non-informative, no structure, no
features, no conditioning. Establishes a floor and a diagnostic: if
naive alone passes playability the task is uninteresting.

\paragraph{Classical ML --- bigram transitions (\texttt{model\_classical.py}).}
Learn three 8$\times$8 matrices, vertical ($P(t_{r,c}\mid t_{r-1,c})$),
horizontal ($P(t_{r,c}\mid t_{r,c-1})$), and a row prior $P(t\mid r)$,
then sample column-by-column. Combine the three via a
\emph{weighted average}
$0.5 \cdot \text{prior} + 0.25 \cdot \text{vert} + 0.25 \cdot \text{horiz}$.
\emph{Why this?} Captures local spatial statistics without the training
burden of a neural model. Key implementation note: I initially tried
multiplicative combination $P \cdot V \cdot H$, which collapsed the
model to all-empty because empty$\to$empty dominates all three
distributions. The weighted average is load-bearing.

\paragraph{Deep learning --- conditional convolutional VAE
(\texttt{model\_deep.py}).}
Encoder: 3 strided conv blocks (8$\to$64$\to$128$\to$256 channels)
with residual connections and BatchNorm, flatten to a 64-dim latent.
Decoder: mirror, conditioned on a 4-dim one-hot (3 difficulty buckets
$+$ 1 null bucket for classifier-free guidance). Input is one-hot 3-class
(collapsed: empty/solid/hazard) on the full 20$\times$40 grid. I use
sqrt-inverse class weights on the cross-entropy reconstruction loss,
$\beta=0.1$ KL weight, and train for 150 epochs on CPU. At inference I
sample $z \sim \mathcal{N}(0, I)$, decode both conditional and
unconditional logits, and extrapolate:
\[
\ell = \ell_\text{uncond} + \gamma (\ell_\text{cond} - \ell_\text{uncond}), \quad \gamma = 2.
\]
\emph{Why this?} Demonstrates a modern deep generative model on the
same task; provides the conditioning signal the UI needs.

\paragraph{Post-processing: \texttt{enforce\_layout} $\to$ \texttt{repair}.}
Every model's raw output is passed through a deterministic post-pass
shared across models: (a) force the sky/floor silhouette, (b) strip
stray sky walls and sub-floor pits, (c) ensure every column has a
floor and a leftmost spawn, and (d) iteratively place one-way
platforms (greedy BFS-guided bridging) and carve walls or hazards
until the right edge is reachable, or 60 iterations elapse. Repair
is the centerpiece of the ablation and the subject of Section~\ref{sec:experiment}.

% ---------------------------------------------------------------------------
\section{Results}
\label{sec:results}

\begin{table}[h]
\centering
\small
\caption{Primary metrics across (model, difficulty) cells on 200
generated samples per cell, with the shipped repair pipeline applied.
\emph{Real@test} is the held-out SuperTux test split as a reference
ceiling. BFS playability is the arc-aware variant described in
Section~\ref{sec:eval}. Arrows indicate direction of improvement.}
\label{tab:main}
\begin{tabular}{lrrrrr}
\toprule
\textbf{Cell} & \textbf{Playability} $\uparrow$ & \textbf{JS div.} $\downarrow$ & \textbf{Ground cov.} & \textbf{Hazard dens.} & \textbf{Diversity} $\uparrow$ \\
\midrule
real@test   & 23.0\% & 0.0111 & 37.4\% & 0.0172 & 0.3475 \\
\midrule
naive@50    & 56.0\% & 0.0193 & 42.9\% & 0.0000 & 0.2796 \\
bigram@50   & \textbf{84.5\%} & \textbf{0.0113} & 35.2\% & 0.0140 & 0.2631 \\
vae@easy    & 31.5\% & 0.0171 & 36.1\% & 0.0604 & 0.2801 \\
vae@medium  & 78.5\% & 0.0114 & 38.3\% & 0.0136 & 0.2180 \\
vae@hard    & 15.5\% & 0.0168 & 31.9\% & 0.0570 & 0.3046 \\
\bottomrule
\end{tabular}
\end{table}

\paragraph{Quantitative comparison.}
Table~\ref{tab:main} summarizes the shipped pipeline (enforce-layout
$\to$ repair) evaluated on 200 samples per cell. Two numbers stand
out:
\begin{itemize}
    \item The bigram exceeds real@test on playability (84.5\% vs.\
    23.0\%). This is not a bug: real@test playability is low because
    SuperTux chunks are slices of longer levels, so many chunks are
    interior screens with no leftmost spawn or no rightmost exit.
    \emph{Repaired} samples are shaped by the repair pass to have
    both, a fair side-effect of the problem framing (I generate
    standalone screens, not slices).
    \item VAE-medium is competitive (78.5\%), but VAE-easy and
    VAE-hard lag. The ablation (Section~\ref{sec:experiment})
    identifies hazard saturation as the root cause on hard and
    ceiling-above-lane geometry on easy.
\end{itemize}

\paragraph{Playable/unplayable outcome matrix.}
Figure~\ref{fig:confusion} is the generative analogue of a confusion
matrix: for each cell, the fraction of 200 samples that are playable
vs.\ unplayable, before and after repair. No-repair bars are near-zero
everywhere; with-repair bars span the full range from VAE-hard (15.5\%)
to bigram (84.5\%).

\begin{figure}[h]
\centering
\includegraphics[width=0.95\linewidth]{repair_ablation.png}
\caption{Playability outcome matrix: 200 samples per (model, difficulty)
cell, before vs.\ after repair. Real@test plotted once as a reference.
Repair yields a +55--85 percentage-point swing on every model cell.}
\label{fig:confusion}
\end{figure}

% ---------------------------------------------------------------------------
\section{Error Analysis}
\label{sec:errors}

I picked five specific mispredictions spanning all five (model,
difficulty) cells, each of which the shipped repair pipeline fails to
make playable. Figure~\ref{fig:errors} shows them; the dashed line
marks the farthest column BFS reaches before halting.

\begin{figure}[h]
\centering
\includegraphics[width=0.9\linewidth]{error_cases.png}
\caption{Five specific mispredictions, one per model cell. Dashed line
shows the column at which BFS halts.}
\label{fig:errors}
\end{figure}

\begin{description}
    \item[Case 1 --- naive@50, seed 4, halts at c=34.] A Gaussian-
    perturbed column of height 14 sits next to one of height 5, and
    the resulting 9-tile step-up exceeds the player's jump reach
    ($\text{max\_jump\_height}=2$, $\text{max\_jump\_width}=2$).
    \emph{Root cause:} the naive baseline has no notion of reachability
    between adjacent columns; height is sampled i.i.d.\ per column.
    \emph{Mitigation:} cap per-column height-delta (a running
    random walk instead of i.i.d.\ Gaussian sampling). One-line fix.

    \item[Case 2 --- bigram@50, seed 7, halts at c=19.] A mid-level
    wall of SOLID tiles rising from the floor to row 4. The bigram's
    row prior says ``column is mostly solid at this row,'' and the
    horizontal transition lets it bleed two columns wide. Repair's
    \texttt{\_carve\_walls} pass can chew through 1--2 tiles but not
    a thick contiguous column.
    \emph{Root cause:} no global reachability awareness in the
    column-by-column sampler.
    \emph{Mitigation:} add a repair pre-pass that thins walls
    $>3$ tiles tall to at most 2 before bridging; or sample entire
    rows conditionally rather than independent columns.

    \item[Case 3 --- vae@easy, seed 0, BFS never enters (no valid spawn).]
    The VAE places hazards on or adjacent to every leftmost standing
    position. Repair's \texttt{\_ensure\_spawn} pass clears a small
    rectangle at the left edge but the surrounding hazards still make
    the spawn unreachable under arc-aware BFS.
    \emph{Root cause:} classifier-free guidance at $\gamma=2$ sharpens
    the easy bucket's tendency to cluster hazards near ground level,
    the conditional signal is ``easy'' but the model has learned
    that visually ``easy'' SuperTux levels include decorative hazards.
    \emph{Mitigation:} enlarge the spawn-clearing rectangle to
    $5 \times 5$, or apply per-bucket hazard-density caps at inference.

    \item[Case 4 --- vae@medium, seed 3, halts at c=34.] A mid-level
    hazard ``moat'' spans four columns at a height the repair can
    theoretically bridge, and indeed repair places a platform above
    it. But the platform height interacts with the arc-clear check:
    the intended jump arc passes through a residual hazard two rows
    above the moat.
    \emph{Root cause:} repair places bridge tiles greedily at a single
    height; it does not iterate on bridge-height choices.
    \emph{Mitigation:} extend \texttt{\_place\_bridge\_step} to try
    multiple heights ($\Delta r \in \{-1, -2, -3\}$) and accept the
    first one whose arc-clear check also passes.

    \item[Case 5 --- vae@hard, seed 1, halts at c=11.] The hard
    bucket's hazard density (6.0\% vs.\ the real data's 1.7\%) saturates
    the level; entire rows are pocked with hazards. Repair can route
    around a small cluster but not a contiguous field.
    \emph{Root cause:} the VAE learns that ``hard'' SuperTux screens
    are hazard-dense, and classifier-free guidance amplifies that
    signal beyond the training distribution.
    \emph{Mitigation:} reduce guidance scale specifically on the hard
    bucket, or cap hazard density at a post-sample step (convert a
    random subset of hazards to empty until density $\leq 0.03$).
\end{description}

Two of the five cases (3 and 5) share a root cause, the VAE's hazard
distribution on the easy and hard buckets is miscalibrated, and a
single mitigation (clip hazard density at inference) would fix both.
That is my top priority under ``Future Work.''

% ---------------------------------------------------------------------------
\section{Experiment: Repair Ablation}
\label{sec:experiment}

\subsection{Experimental plan}
The research question is: \emph{how much of the shipped app's
playability comes from the model, and how much from the BFS repair
post-pass?} I hold everything else constant and vary only the
presence of \texttt{scripts.repair.repair}. For each of six cells
(real@test, naive@50, bigram@50, vae@easy, vae@medium, vae@hard) I
generate 200 levels with fixed seed 42 and compute all five metrics
from Section~\ref{sec:eval} under two conditions:
\begin{enumerate}
    \item \textbf{no\_repair}: model sample $\to$
    \texttt{enforce\_layout} $\to$ 3-class collapse $\to$ metrics.
    \item \textbf{with\_repair}: model sample $\to$
    \texttt{enforce\_layout} $\to$ \texttt{repair} $\to$ 3-class
    collapse $\to$ metrics.
\end{enumerate}
Both conditions share \texttt{enforce\_layout}, so the only delta is
the BFS repair loop itself.

\subsection{Results}

\begin{table}[h]
\centering
\small
\caption{Repair ablation, 200 samples per cell, seed=42. $\Delta$ is
with\_repair $-$ no\_repair.}
\label{tab:ablation}
\begin{tabular}{lrrr|rrr}
\toprule
 & \multicolumn{3}{c|}{\textbf{Playability} $\uparrow$} & \multicolumn{3}{c}{\textbf{JS divergence} $\downarrow$} \\
\textbf{Cell} & no-rep. & with-rep. & $\Delta$ & no-rep. & with-rep. & $\Delta$ \\
\midrule
naive@50    & 0.0\% & 56.0\% & \textbf{+56.0} & 0.0225 & 0.0193 & $-0.0032$ \\
bigram@50   & 0.0\% & 84.5\% & \textbf{+84.5} & 0.0112 & 0.0113 & $+0.0001$ \\
vae@easy    & 0.0\% & 31.5\% & \textbf{+31.5} & 0.0191 & 0.0171 & $-0.0020$ \\
vae@medium  & 0.5\% & 78.5\% & \textbf{+78.0} & 0.0119 & 0.0114 & $-0.0004$ \\
vae@hard    & 0.0\% & 15.5\% & \textbf{+15.5} & 0.0172 & 0.0168 & $-0.0004$ \\
\bottomrule
\end{tabular}
\end{table}

\subsection{Interpretation}
Three findings.

\paragraph{Finding 1: repair is where the playability comes from.}
Raw model output is 0\% playable on four of five cells (0.5\% on
vae@medium). Every shipped beatable level is a repaired level. This
inverts the intuition that the model ``learns'' the constraint --- on
a dataset this small, with chunks that are already sub-slices of
longer levels, the generative model has no incentive to output
left-to-right traversable geometry. The repair pass encodes that
constraint explicitly.

\paragraph{Finding 2: repair is near-free on distribution.}
JS divergence shifts by at most 0.0032 and usually less than 0.001 ---
an order of magnitude below inter-model variance. Repair does not
distort the tile histogram into a ``repaired'' distribution; it
surgically flips a small number of tiles near the BFS cut.

\paragraph{Finding 3: repair has failure modes that point at the model.}
Bigram's +84.5 lift vs.\ VAE-hard's +15.5 lift identifies where the
model's output exceeds repair's capacity. Hazard-dense VAE output
(hard bucket, 6.0\% hazard density) defeats the repair pass because
repair is reachability-focused: it adds platforms, it does not remove
hazards in bulk. This is the clearest signal that the next investment
should be in reshaping the VAE's class distribution (or hazard-capping
at inference), not in extending the repair logic.

\subsection{Recommendations}
\begin{itemize}
    \item Keep repair as the primary playability enforcer. It is
    cheap, deterministic, and easy to debug (unlike a learned
    constraint).
    \item Address the VAE hazard miscalibration directly: clip hazard
    density per bucket at inference.
    \item Do \emph{not} bet on training the playability constraint
    into the model at this data scale. The ablation shows the model
    lacks the signal.
\end{itemize}

% ---------------------------------------------------------------------------
\section{Conclusions}

Three generative models produce visually plausible but unplayable
output; a shared BFS-guided repair pass makes them shippable. The
bigram transition model, helped by repair, reaches 84.5\% playability,
higher than the conditional VAE does at any difficulty. The
take-away for short-horizon tile-grid PCG is that hard constraints
(playability) are better handled outside the learned model than
inside it, at least at this data scale.

The system is deployed at \url{https://arnav161-genterrain.hf.space};
users can pick a model, a difficulty, and either play a single
generated level or play an endless seam-stitched stream of chunks.

% ---------------------------------------------------------------------------
\section{Future Work}

\begin{itemize}
    \item \textbf{Hazard-calibrated VAE sampling.} Implement
    inference-time hazard capping (Section~\ref{sec:errors}); expected
    to lift VAE-easy and VAE-hard playability substantially.
    \item \textbf{Larger dataset via SuperTux add-ons.} The community
    add-on repo hosts $\sim$10$\times$ more level data; rerunning the
    full pipeline on that would let me compare naive/bigram/VAE at a
    more representative scale. Part of the reason repair dominates
    here is data starvation.
    \item \textbf{Sequence model for long horizons.} For endless mode
    I currently stitch 20$\times$40 chunks by copying the rightmost
    column as context for the next chunk. A small Transformer over
    column-sequences would give the model actual long-range context
    (and remove the stitching-discontinuity failure modes I see in
    endless mode today).
    \item \textbf{Human playtesting.} BFS playability is a proxy for
    ``fun.'' A paired user study (\emph{which level feels more fun?})
    would tell me whether VAE-diversity or bigram-consistency matters
    more to actual players.
    \item \textbf{Broader experiment grid.} The repair ablation was a
    single axis. A 3-axis factorial (vocabulary size $\times$
    repair-iterations $\times$ guidance scale) would surface
    higher-order effects.
\end{itemize}

% ---------------------------------------------------------------------------
\section{Commercial Viability Statement}

\paragraph{Where this works.} The pipeline is a reasonable starting
point for (a) indie 2D platformer studios wanting a
randomized-daily-level mode with a hard playability guarantee, or
(b) educational games that need a steady stream of short, beatable
puzzle screens. The BFS repair pass is strong enough that every
shipped level is beatable; the content is derivative (SuperTux-style
sidescrollers) but the infrastructure (repair, seam-stitching,
metrics) is re-usable over any tile-grid game with a similar
left-to-right traversal model.

\paragraph{Where it does not work yet.} Three gaps prevent direct
commercial use:
\begin{enumerate}
    \item \textbf{Aesthetic ceiling.} Output is playable but not
    interesting. A diversity score of 0.26 means many levels read
    similarly. Commercial players would churn.
    \item \textbf{Single-game training data.} SuperTux's 1{,}500 chunks
    are not enough to generalize out of style. A commercial deployment
    would need a much larger, licensable content dataset.
    \item \textbf{No game-design signal.} The model doesn't know
    anything about pacing, difficulty progression, or tutorial
    design. Real platformer levels are curated sequences, not
    independent screens.
\end{enumerate}

This is a useful infrastructure
demonstration and a valid research project, but it is not a product.
A commercial version would inherit the repair pass and the metric
framework, and replace the generative model with a larger
transformer-based one trained on licensed content.

% ---------------------------------------------------------------------------
\section{Ethics Statement}

\paragraph{Training-data provenance.}
SuperTux is MIT-licensed; re-using the levels for model training and
rendering generated content is permissible under that license. I do
not redistribute the SuperTux assets themselves in the repository;
only the derived category mapping and processed chunk arrays are
shipped.

\paragraph{Content risks.}
The model generates tile grids. It does not produce text, images of
people, or any content with plausible routes to discriminatory,
explicit, or privacy-violating output. The main content risk is
\emph{derivative style}: a generated level is trivially
recognizable as SuperTux-adjacent. I explicitly do not claim
originality of theme.

\paragraph{Creator displacement.}
If systems like this were to scale to a large, licensed training set
and be deployed as a replacement for human level designers, the
standard generative-AI-economics concerns apply: level-designer work
is a creative labor market, and automating it absent a licensing and
compensation structure would harm designers. This project is
deliberately small-scale and sourced from an open-source community;
any commercial extension should pay the source creators.

\paragraph{Gameplay fairness.}
The \emph{endless} mode posts scores to a leaderboard. Because the
generator is stochastic, different players may see different levels
of difficulty on different runs. But a strict competitive mode would need
a shared-seed option for fairness.

\paragraph{User data.}
The deployed app stores only a username and hashed password per user
account, plus opt-in gameplay stats (completions by difficulty,
endless-mode best). No personal information beyond the self-chosen username; guest
mode requires no account.

% ---------------------------------------------------------------------------
\section*{Reproducibility}

Code, data pipeline, and trained models are at\\
\url{https://github.com/arnavmahale/game-level-generation}.
All numbers in this report come from:
\begin{itemize}
    \item \texttt{scripts/experiment\_repair\_ablation.py\\ --n 200 --seed 42}
    (Tables~\ref{tab:main} and~\ref{tab:ablation}, Figure~\ref{fig:confusion}).
    \item \texttt{scripts/report\_error\_cases.py} (Figure~\ref{fig:errors}).
\end{itemize}
Both scripts run on CPU in under five minutes from the trained-model
state.

% ---------------------------------------------------------------------------
\bibliographystyle{plainnat}
\bibliography{references}

\end{document}
